\chapter*{Introduzione}
\chaptermark{Introduzione}
I neutrini ($\nu$) si collocano all'interno del Modello Standard delle particelle elementari nella famiglia dei leptoni inizialmente descritte come prive di massa. A differenza degli altri leptoni, il neutrino interagisce solo tramite l'interazione debole, poichè elettricamente neutro. Esiste in tre sapori: elettronico ($\nu_e$), muonico ($\nu_\mu$) e tauonico ($\nu_\tau$), ciascuno associato al rispettivo leptone con cui viene generato o assorbito durante un'interazione.

Dalla prima ipotesi formulata da W. Pauli nel 1930 alla loro rivelazione sperimentale nel 1956, ad opera di Clyde Cowan e Frederick Reines, lo studio dei neutrini ha attraversato numerose fasi. In particolare, il problema dei neutrini solari ha messo in luce una discrepanza tra le previsioni teoriche e le osservazioni sperimentali. Successivamente, esperimenti di grande rilevanza come SNO (Sudbury Neutrino Observatory) e Super-KamiokaNDE hanno confermato l'esistenza di un fenomeno quantistico legato alla natura di queste particelle: l’\textbf{oscillazione dei neutrini}.

Dalla seconda metà del Novecento sono stati compiuti enormi progressi nella comprensione della fenomenologia dei neutrini, malgrado le numerose difficoltà legate alla loro rivelazione. Nonostante ciò, attualmente rimangono diverse domande aperte: da un lato, si cerca di affinare la misura dei parametri già noti; dall’altro, si tenta di esplorare proprietà ancora sconosciute, come l’ordinamento di massa dei neutrini (NMO). In questo contesto si colloca il rivelatore JUNO (Jiangmen Underground Neutrino Observatory), che rappresenta una delle migliori opportunità per fare luce su questi aspetti ancora irrisolti.

JUNO è un rivelatore situato in Cina, a $52\,\unit{km}$ da due centrali nucleari. Il nucleo dell'esperimento è costituito da $20\,\unit{kton}$ di scintillatore liquido contenuto nel \textbf{Central Detector} (CD), una sfera dal diametro di $35.4\,\unit{m}$. Quest'ultima è inserita all'interno della \textbf{Water Pool} (WP), un cilindro che contiene $35\,\unit{kton}$ di acqua ultra pura.   Grazie alle sue dimensioni e alla sua eccellente risoluzione energetica, si prevede che JUNO possa ottenere risultati di grande rilievo. Infatti, misurando i flussi di $\bar{\nu}_e$ emessi dai reattori nucleari, sarà in grado di fornire un contributo significativo nella determinazione dell’ordinamento di massa dei neutrini, con una significatività statistica di $3\sigma$ in sette anni di raccolta dati. Oltre a questo, JUNO contribuirà a numerosi altri obiettivi scientifici dalla misura dei \textbf{geoneutrini}, dei \textbf{neutrini atmosferici} e dei \textbf{neutrini solari}, oltre alla rivelazione dei \textbf{neutrini da supernova}.

Un aspetto essenziale per la maggior parte degli obiettivi scientifici di JUNO riguarda l'ottimizzazione della riduzione dei segnali di background. In particolare, per alcune misure di neutrini solari, un contributo significativo al segnale è dato dal background cosmogenico, dovuto al decadimento di isotopi radioattivi prodotti dai muoni cosmici che attraversano il CD e interagiscono con lo scintillatore liquido. Risulta quindi essenziale l'implementazione di tecniche di veto, sviluppate per escludere questi eventi dalle misure. Una di queste tecniche si basa sulla classificazione e ricostruzione precisa delle tracce dei \textbf{muoni}.

Attualmente, la costruzione di JUNO è completata e il rivelatore si trova nella \textbf{fase di riempimento} con lo scintillatore nel Central Detector. Durante questa fase, è possibile caratterizzare il flusso di muoni grazie all’abbondante produzione di \textbf{luce Cherenkov}. I muoni che attraverano l'intero rivelatore hanno in media un’energia di circa $209\,\unit{GeV}$ e generano un segnale luminoso ben identificabile. Di conseguenza, l’analisi dei muoni non dipende  dallo stato di riempimento con lo scintillatore, permettendo di raccogliere dati utili sin dalle prime fasi operative.
Il presente lavoro di tesi si concentra sull’analisi dei primi mesi di raccolta dati di JUNO, da gennaio 2025 a marzo 2025, con l'obiettivo di caratterizzare il flusso e le proprietà distintive degli eventi identificati come muoni.

Il presente elaborato è articolato come segue:
\begin{itemize}
    \item nel primo capitolo viene presentata una panoramica generale sulla fisica dei neutrini, con particolare attenzione agli esperimenti che hanno caratterizzato il "problema dei neutrini solari";
    \item nel secondo capitolo viene descritto l'esperimento JUNO, illustrandone gli obiettivi scientifici e lo stato attuale del rivelatore;
    \item nel terzo capitolo viene presentato lo studio sulla caratterizzazione del flusso di muoni, descrivendo i dati analizzati, i due metodi utilizzati per l'analisi e i risultati ottenuti.
\end{itemize}
